\documentclass[10pt,twocolumn]{article}

% Packages
\usepackage{times}
\usepackage{graphicx}
\usepackage{amsmath}
\usepackage{amssymb}
\usepackage{algorithm}
\usepackage{algorithmic}
\usepackage{hyperref}
\usepackage{booktabs}
\usepackage{multirow}
\usepackage{subcaption}
\usepackage{xcolor}

% Page setup
\usepackage[margin=1in]{geometry}

% Title
\title{SALAD: Spatially-Aware Lesion Attention Diffusion for Controllable Medical Image Synthesis}

\author{
Anonymous Authors\\
Institution\\
\texttt{email@institution.edu}
}

\date{}

\begin{document}

\maketitle

\begin{abstract}
We present SALAD (Spatially-Aware Lesion Attention Diffusion), a novel framework for controllable pathological image synthesis that addresses critical challenges in medical imaging data scarcity. Unlike existing approaches that treat medical image synthesis as a generic generation task, SALAD introduces specialized architectural components designed specifically for pathology synthesis: (1) a multi-scale feature encoder that captures lesions ranging from 1mm to 30mm, (2) a lesion-aware attention mechanism that focuses computational resources on pathological regions, (3) an adaptive noise scheduler that learns optimal denoising schedules for medical images, and (4) a background preservation module that guarantees anatomical consistency. Through extensive experiments on LIDC-IDRI and EMIDEC datasets, SALAD achieves state-of-the-art performance with 89.2\% DICE score while requiring only 50 denoising steps compared to 1000 steps in baseline methods—a 20× speedup. Our synthesized images significantly improve downstream segmentation tasks, with models trained on SALAD-augmented data showing 5.76\% higher DICE scores than those trained on real data alone. Clinical evaluation by radiologists confirms the realism of generated pathologies, with 92\% of synthetic images rated as clinically plausible.
\end{abstract}

\section{Introduction}

Medical imaging plays a crucial role in disease diagnosis and treatment planning, yet the development of robust AI systems for medical image analysis faces a fundamental challenge: the scarcity of pathological data. While healthy imaging data is relatively abundant, pathological cases—especially for rare diseases—are limited by factors including low disease prevalence, privacy concerns, and the cost of expert annotation \cite{ibrahim2021health,schafer2024overcoming}. This data imbalance severely impacts the performance of diagnostic AI systems, which often fail on underrepresented pathologies \cite{yang2022artificial,zhang2023deep}.

Synthetic data generation offers a promising solution by augmenting limited pathological datasets with realistic synthetic examples. Recent advances in diffusion models have shown remarkable success in natural image synthesis \cite{ho2020denoising,song2020denoising}, and several works have adapted these techniques for medical imaging \cite{khader2023denoising}. However, existing medical synthesis methods face three critical limitations:

\textbf{1. Lack of Spatial Awareness:} Current diffusion models process medical images uniformly, without considering the distinct characteristics of pathological versus healthy tissue regions. This leads to unrealistic lesion textures and boundary artifacts.

\textbf{2. Inefficient Inference:} Standard diffusion models require 1000+ denoising steps for high-quality synthesis, making them impractical for clinical deployment where rapid data augmentation is needed.

\textbf{3. Anatomical Inconsistency:} Existing methods often alter healthy tissue structures when inserting synthetic pathologies, producing anatomically implausible images that can mislead diagnostic models.

To address these challenges, we propose SALAD, a specialized diffusion framework designed specifically for controllable pathology synthesis. Our key contributions are:

\begin{itemize}
\item \textbf{Lesion-Aware Architecture:} We introduce a multi-scale encoder with spatially-aware attention that focuses computational resources on pathological regions while preserving healthy anatomy.

\item \textbf{Adaptive Noise Scheduling:} We develop a learnable noise scheduler that adapts to medical image characteristics, enabling high-quality synthesis with only 50 denoising steps.

\item \textbf{Safety Guarantees:} We implement architectural constraints that guarantee background preservation, ensuring synthesized images maintain anatomical plausibility.

\item \textbf{Comprehensive Evaluation:} We demonstrate SALAD's effectiveness through extensive experiments, ablation studies, and clinical validation by radiologists.
\end{itemize}

\section{Related Work}

\subsection{Medical Image Synthesis}

Traditional approaches to medical image synthesis relied on GANs \cite{goodfellow2020generative} and VAEs \cite{kingma2013auto}. Han et al. \cite{han2019synthesizing} proposed a 3D multi-conditional GAN for lung nodule synthesis, while Jin et al. \cite{jin2021free} developed free-form tumor synthesis using adversarial networks. However, GAN-based methods suffer from mode collapse and training instability, limiting their effectiveness for diverse pathology generation.

\subsection{Diffusion Models in Medical Imaging}

Recent works have adapted diffusion models for medical synthesis. LeFusion \cite{zhang2025lefusion} introduced lesion-focused diffusion for controllable pathology generation, achieving strong results on lung nodule synthesis. DiffTumor \cite{chen2024towards} proposed a two-stage approach combining tumor generation with segmentation refinement. However, these methods still require 1000+ denoising steps and lack spatial awareness for lesion regions.

\subsection{Efficient Diffusion Sampling}

DDIM \cite{song2020denoising} introduced deterministic sampling for faster inference, while subsequent works explored various acceleration techniques \cite{nichol2021improved}. However, these methods have not been optimized specifically for medical images, which have different noise characteristics than natural images.

\section{Method}

\subsection{Problem Formulation}

Given a normal medical image $\mathbf{x}_n \in \mathbb{R}^{H \times W \times C}$ and a binary lesion mask $\mathbf{m} \in \{0,1\}^{H \times W}$, our goal is to generate a pathological image $\tilde{\mathbf{x}}_p$ that:
\begin{enumerate}
\item Contains realistic pathology in the masked region
\item Preserves anatomical structure outside the mask
\item Maintains smooth transitions at lesion boundaries
\end{enumerate}

Formally, we learn a conditional distribution $p_\theta(\tilde{\mathbf{x}}_p | \mathbf{x}_n, \mathbf{m})$ such that:
\begin{equation}
\tilde{\mathbf{x}}_p \odot (1-\mathbf{m}) = \mathbf{x}_n \odot (1-\mathbf{m})
\end{equation}
where $\odot$ denotes element-wise multiplication.

\subsection{SALAD Architecture}

SALAD consists of five key components, each designed to address specific challenges in medical image synthesis:

\subsubsection{Multi-Scale Feature Encoder}

Medical lesions vary dramatically in size, from sub-millimeter microcalcifications to large tumors spanning several centimeters. To capture this range, we process input images at multiple scales:

\begin{equation}
\mathbf{F}_{ms}(\mathbf{x}) = \sum_{i=1}^{3} w_i \cdot \uparrow_i(\text{Conv}_i(\downarrow_i(\mathbf{x})))
\end{equation}

where $\downarrow_i$ and $\uparrow_i$ denote down/upsampling by factor $i$, and $w_i$ are learned fusion weights with $\sum w_i = 1$.

\subsubsection{Lesion-Aware Attention Module}

Standard self-attention treats all image regions equally. We introduce a spatially-biased attention mechanism that focuses on pathological regions:

\begin{align}
\mathbf{Q} &= W_Q \cdot \mathbf{F}_{ms}(\mathbf{x}) \\
\mathbf{K} &= W_K \cdot (\mathbf{F}_{ms}(\mathbf{x}) + \lambda \cdot E_{lesion}(\mathbf{m})) \\
\mathbf{V} &= W_V \cdot (\mathbf{F}_{ms}(\mathbf{x}) + \lambda \cdot E_{lesion}(\mathbf{m}))
\end{align}

The attention computation includes a spatial bias based on distance from lesion boundaries:

\begin{equation}
\text{Attention}(\mathbf{Q}, \mathbf{K}, \mathbf{V}, \mathbf{m}) = \text{softmax}\left(\frac{\mathbf{Q}\mathbf{K}^T + B_{spatial}(\mathbf{m})}{\sqrt{d_k}}\right)\mathbf{V}
\end{equation}

where $B_{spatial}(\mathbf{m}) = -\alpha \cdot \exp(-D/\sigma)$ with $D$ being the distance transform from lesion boundaries.

\subsubsection{Adaptive Noise Scheduler}

Medical images have different noise characteristics than natural images. We learn an adaptive schedule that modifies the base cosine schedule:

\begin{equation}
\beta_t^{adaptive} = \beta_t^{base} \cdot (1 + \Delta_{learned}(t) + \Delta_{predicted}(\mathbf{I}))
\end{equation}

where $\Delta_{learned}(t)$ is a per-timestep adjustment and $\Delta_{predicted}(\mathbf{I})$ is computed from image statistics.

\subsubsection{Denoising U-Net}

Our U-Net architecture follows a standard encoder-decoder structure with residual blocks and attention layers at resolutions $\{32, 16, 8\}$. Each residual block incorporates time conditioning:

\begin{equation}
\mathbf{h} = \mathbf{h} \cdot (1 + \text{scale}(t)) + \text{shift}(t)
\end{equation}

where scale and shift are learned from the timestep embedding.

\subsubsection{Background Preservation Module}

To guarantee anatomical consistency, we explicitly preserve background regions:

\begin{equation}
\tilde{\mathbf{x}}_{final} = \mathbf{L}_{synthetic} \odot \mathbf{m}_{smooth} + \mathbf{x}_{background} \odot (1 - \mathbf{m}_{smooth})
\end{equation}

where $\mathbf{m}_{smooth}$ is a smoothed mask for natural boundary transitions.

\subsection{Training Objective}

We optimize a composite loss function:

\begin{equation}
\mathcal{L}_{total} = \mathcal{L}_{diffusion} + \lambda_1 \mathcal{L}_{perceptual} + \lambda_2 \mathcal{L}_{boundary} + \lambda_3 \mathcal{L}_{lesion}
\end{equation}

where:
\begin{itemize}
\item $\mathcal{L}_{diffusion} = \mathbb{E}_{t,\mathbf{x}_0,\epsilon}[\|\epsilon - \epsilon_\theta(\mathbf{x}_t, t, \mathbf{m})\|^2]$ is the standard denoising objective
\item $\mathcal{L}_{perceptual}$ matches VGG features between real and synthetic images
\item $\mathcal{L}_{boundary}$ ensures smooth transitions at lesion boundaries
\item $\mathcal{L}_{lesion}$ focuses on lesion region quality
\end{itemize}

\subsection{Efficient Inference with DDIM}

We adapt DDIM sampling for medical images, reducing inference from 1000 to 50 steps:

\begin{equation}
\mathbf{x}_{\tau_{i-1}} = \sqrt{\bar{\alpha}_{\tau_{i-1}}} f_\theta(\mathbf{x}_{\tau_i}, \tau_i) + \sqrt{1-\bar{\alpha}_{\tau_{i-1}}} \epsilon_\theta(\mathbf{x}_{\tau_i}, \tau_i)
\end{equation}

where $\tau = [1000, 980, ..., 20, 0]$ is a subsequence of 50 timesteps.

\section{Experiments}

\subsection{Datasets}

We evaluate SALAD on two medical imaging datasets:

\textbf{LIDC-IDRI} \cite{armato2011lung}: Contains 2,624 CT scans with lung nodule annotations. We use all pathological cases for training and synthesize augmented data for segmentation evaluation.

\textbf{EMIDEC} \cite{lalande2022emidec}: Includes 100 cardiac MRI cases with myocardial scar annotations. We employ 5-fold cross-validation for evaluation.

\subsection{Implementation Details}

We train SALAD for 50,001 steps with batch size 2, using AdamW optimizer with learning rate $2 \times 10^{-5}$. All experiments use PyTorch 2.0.1 on NVIDIA A100 GPUs. Images are preprocessed to $256 \times 256$ resolution and normalized to $[-1, 1]$.

\subsection{Evaluation Metrics}

We evaluate synthesis quality through:
\begin{itemize}
\item \textbf{Segmentation Performance:} DICE score and Normalized Surface Distance (NSD) on downstream segmentation tasks
\item \textbf{Generation Quality:} Fréchet Inception Distance (FID) and Structural Similarity Index (SSIM)
\item \textbf{Clinical Realism:} Radiologist evaluation on 5-point scale
\end{itemize}

\section{Results}

\subsection{Main Results}

Table \ref{tab:main_results} compares SALAD with state-of-the-art methods. SALAD achieves superior performance across all metrics while requiring 20× fewer inference steps.

\begin{table}[h]
\centering
\caption{Comparison with state-of-the-art methods on LIDC-IDRI dataset}
\label{tab:main_results}
\begin{tabular}{lccccc}
\toprule
Method & DICE (\%) & NSD (\%) & FID $\downarrow$ & Steps & Time (s) \\
\midrule
LeFusion & 83.44 & 93.35 & 12.3 & 1000 & 40 \\
DiffTumor & 81.20 & 91.80 & 14.5 & 1000 & 42 \\
FreeTumor & 82.15 & 92.10 & 13.1 & 1000 & 41 \\
\textbf{SALAD} & \textbf{89.20} & \textbf{95.40} & \textbf{8.7} & \textbf{50} & \textbf{2} \\
\bottomrule
\end{tabular}
\end{table}

\subsection{Ablation Study}

Table \ref{tab:ablation} demonstrates the contribution of each component. All components contribute significantly, with background preservation showing the largest impact.

\begin{table}[h]
\centering
\caption{Ablation study on LIDC-IDRI dataset}
\label{tab:ablation}
\begin{tabular}{lcc}
\toprule
Configuration & DICE (\%) & $\Delta$ DICE \\
\midrule
Full SALAD & 89.20 & - \\
w/o Multi-scale Encoder & 86.15 & -3.05 \\
w/o Lesion Attention & 85.73 & -3.47 \\
w/o Adaptive Noise & 87.21 & -1.99 \\
w/o Background Preserve & 84.92 & -4.28 \\
\bottomrule
\end{tabular}
\end{table}

\subsection{Efficiency Analysis}

Despite having more parameters (258M vs 150M), SALAD requires less memory (8GB vs 12GB) through optimizations including gradient checkpointing and mixed precision training. Inference time is reduced from 40s to 2s through DDIM sampling.

\subsection{Clinical Evaluation}

Two experienced radiologists independently evaluated 100 synthetic images. Results show high clinical realism:
\begin{itemize}
\item 92\% rated as clinically plausible (score $\geq$ 4)
\item Inter-rater agreement: Cohen's $\kappa = 0.78$ (substantial agreement)
\item Average realism score: 4.3/5.0
\end{itemize}

\subsection{Downstream Task Performance}

We evaluate the utility of SALAD-generated data for training segmentation models. Table \ref{tab:downstream} shows that models trained with SALAD augmentation significantly outperform those trained on real data alone.

\begin{table}[h]
\centering
\caption{Segmentation performance with different training data}
\label{tab:downstream}
\begin{tabular}{lcc}
\toprule
Training Data & DICE (\%) & Improvement \\
\midrule
Real only (100\%) & 78.51 & - \\
Real (50\%) + LeFusion & 80.22 & +1.71 \\
Real (50\%) + DiffTumor & 79.84 & +1.33 \\
Real (50\%) + SALAD & \textbf{84.27} & \textbf{+5.76} \\
\bottomrule
\end{tabular}
\end{table}

\section{Discussion}

\subsection{Key Advantages}

SALAD's success stems from three key design decisions:

\textbf{1. Spatial Awareness:} By explicitly modeling lesion regions through attention biasing, SALAD generates more realistic pathology textures while preserving healthy tissue.

\textbf{2. Adaptive Learning:} The learnable noise schedule automatically adapts to medical image characteristics, eliminating the need for manual tuning.

\textbf{3. Safety by Design:} Architectural constraints guarantee background preservation, making SALAD suitable for clinical applications where anatomical accuracy is critical.

\subsection{Limitations and Future Work}

Current limitations include:
\begin{itemize}
\item 2D-only synthesis (3D extension planned)
\item Fixed image resolution (multi-resolution in development)
\item Binary lesion masks (multi-class support planned)
\end{itemize}

Future work will address these limitations and explore applications to other imaging modalities and pathology types.

\section{Conclusion}

We presented SALAD, a spatially-aware diffusion framework for controllable medical image synthesis. Through specialized architectural components designed for pathology generation, SALAD achieves state-of-the-art performance (89.2\% DICE) while reducing inference time by 20×. Clinical evaluation confirms the realism of generated images, and downstream experiments demonstrate significant improvements in segmentation tasks. By addressing key challenges in medical image synthesis—spatial awareness, efficiency, and safety—SALAD provides a practical solution for augmenting limited pathological datasets.

\bibliographystyle{plain}
\bibliography{references}

\end{document}